\markdownRendererDocumentBegin
\markdownRendererUlBegin
\markdownRendererUlItem change point analysis for each region?\markdownRendererInterblockSeparator
{}\markdownRendererUlBeginTight
\markdownRendererUlItem does a stat sig. chng pt negate a stat sig trend?\markdownRendererUlItemEnd 
\markdownRendererUlEndTight \markdownRendererUlItemEnd 
\markdownRendererUlItem should we subset recent data to match frequency of older data? or should we repeate analysis using the daily interpolated data?\markdownRendererUlItemEnd 
\markdownRendererUlItem include subsurface temp data from moorings\markdownRendererInterblockSeparator
{}\markdownRendererUlBeginTight
\markdownRendererUlItem calculate correlations with ice metrics\markdownRendererUlItemEnd 
\markdownRendererUlItem turnover dates correlations?\markdownRendererUlItemEnd 
\markdownRendererUlEndTight \markdownRendererUlItemEnd 
\markdownRendererUlItem mooring data might be doable but limited regions and years\markdownRendererUlItemEnd 
\markdownRendererUlItem calculate first last and ndays using the ACTUAL eqn you show\markdownRendererUlItemEnd 
\markdownRendererUlItem also calculate the DAY OF MAX\markdownRendererUlItemEnd 
\markdownRendererUlItem Theil-Sen estimates\markdownRendererUlItemEnd 
\markdownRendererUlItem read theil-sen paper (siegel 1982)\markdownRendererUlItemEnd 
\markdownRendererUlItem confidence intervals (how to show on plot?)\markdownRendererUlItemEnd 
\markdownRendererUlItem is some form of trend analysis possible using a subset of the data?\markdownRendererInterblockSeparator
{}\markdownRendererUlBeginTight
\markdownRendererUlItem specifically, if we categorize high/med/low ice years can we perform linear regresssion on each set?\markdownRendererUlItemEnd 
\markdownRendererUlItem seems like discontinuity (i.e. varying delta t) would be a problem\markdownRendererUlItemEnd 
\markdownRendererUlEndTight \markdownRendererUlItemEnd 
\markdownRendererUlItem reconcile number of subregions on all plots\markdownRendererInterblockSeparator
{}\markdownRendererUlBeginTight
\markdownRendererUlItem shouldn't SMR be part of huron?;\markdownRendererUlItemEnd 
\markdownRendererUlEndTight \markdownRendererUlItemEnd 
\markdownRendererUlEnd \markdownRendererDocumentEnd